\section{Future Additions}\label{future}

In this section of the document, we discuss future additions to the Blossom protocol which we have not yet implemented. However, these additions may provide additional performance increases from those previously displayed. The additional functionalities we mention in this section are not guaranteed to be implemented or researched by our team. Rather, these implementations are highlighted due to their unique technologies, which when proven, may provide meaningful benefits to the Blossom protocol. The specific approaches are as follows:

\subsection{Aggregated Signature Schemes}
Aggregated signature schemes would allow for multiple signatures to be expressed as a single signature. Implemented in Blossom such an algorithm would allow all of the transactions stored within a single block to be validated by use of a single signature. This would drastically reduce the computational overhead of signature verification by several orders of magnitude (dependent on the number of transactions in a block). While this technology is promising, the algorithms are still in the research phase with other blockchain foundations working on implementing the technology within their protocols. However, despite the successful implementation of this novel signature scheme, there are serious security worries which may discredit this approach. As to not speculate we will avoid discussing the possible tradeoffs of this technology in detail; however, we would certify that we would not implement a novel signature scheme unless it was guaranteed to have a statistically impossible likelihood of being broken. In the meantime, we are focusing on research on advancing the efficiency of Blossom with existing signature schemes, and it is our belief that given any possible future innovation, the work we are doing in the present to extend the functionality of our protocol is a much more valuable endeavor than speculating on unproven technology.

\subsection{Quantum Safe Cryptography} 
Quantum Safe (QS) crypto is an overarching classification for novel crypto algorithms which are intended to resist attacks from classical computers, as well as quantum computers. Over the past six (6) years the National Institute of Standards and Technology (NIST) has undertaken a process of validating existing QS algorithms, with the selected winners undergoing a standardization process \cite{Boutin2022-jl}. We have been following the developments of the NIST Post-Quantum Cryptography (PQC) competition, and following the 2022 announcement of the four (4) preliminary winners, we look forward to observing the results of those algorithms following the multi-year standardization process. The present algorithms, while promising, still process many trade-offs which would make their implementation extremely difficult for Blossom due to the high computational cost of verifying QS signatures. However, given the immense amount of progress in the field over the last decade, we expect more performant algorithms will be developed in time for the coming proliferation of quantum computers.
